\section{Introduction}
\label{sec:intro}

Self-reflection is becoming a core building block of LLM agent architectures. Systems like Self-Refine \cite{madaan2023selfrefine} and Reflexion \cite{shinn2023reflexion} prompt models to critique their own outputs and revise accordingly, with the expectation that this loop moves the model closer to correct answers. Yet a growing body of evidence suggests the opposite: \citet{huang2024selfcorrect} showed that intrinsic self-correction---without external feedback---consistently \emph{degrades} performance across multiple benchmarks and models. This creates a puzzle. If self-critique changes outputs (sometimes dramatically), what is happening inside the model? Does the critique-revision cycle restructure the model's internal representations, or does it merely re-sample from the same distribution with a perturbed prompt?

This question matters because internal representations are the substrate of model computation. Prior work has shown that residual stream activations encode rich information about reasoning quality: models internally represent correct answers even when they output wrong ones \cite{garner2025probing}, and linear probes can predict eventual chain-of-thought correctness from early hidden states \cite{david2025temporal}. If self-critique genuinely restructures these representations, it could in principle move them toward regions that encode better reasoning. If instead it merely perturbs them, the observed behavioral degradation would follow naturally.

We address this gap by directly measuring how residual stream activations shift across the generate-critique-revise pipeline. Using \qwen on 200 \gsmk problems, we extract activations from all 28 transformer layers at four stages: initial answer, self-critique, revised answer, and a control condition (re-prompting without critique). We then quantify drift via cosine similarity, train linear probes for correctness prediction, and analyze subspace geometry via PCA.

Our key finding is that self-critique induces significantly more representational drift than re-prompting ($p < 0.01$ at all 28 layers after Benjamini-Hochberg correction), but this drift is destructive---accuracy drops by 6.5 percentage points, linear probe separability decreases, and representations move away from correct reasoning subspaces. Self-critique is not shallow re-sampling; it is genuine restructuring in the wrong direction.

We make the following contributions:
\begin{itemize}[leftmargin=*,itemsep=0pt,topsep=0pt]
    \item We provide the first study tracking residual stream representation drift across the self-critique pipeline, showing that critique induces significantly larger shifts than prompt-variation controls at all transformer layers.
    \item We demonstrate that this representational restructuring is destructive: it degrades linear separability of correctness and moves activations away from initially well-calibrated regions.
    \item We connect behavioral observations (accuracy degradation) to mechanistic evidence (representation drift), bridging the gap between prior work on self-correction failure and probing-based interpretability.
    \item We identify layer-specific drift patterns, finding that middle layers (7--14) are most affected by self-critique, with implications for targeted intervention strategies.
\end{itemize}
